We identified a sample of 166 accelerating stars that may host substellar companions and do not have archival spectra. From APO/ARCES spectra of the full 166 star sample we measured lithium equivalent widths (or upper limits) and $R^{'}_{HK}$. We acquired TESS light curves for 129 stars and measured rotation rates for 23 out of the 129 stars. For the earlier spectral types, we also constrain the ages from color-magnitude diagram positions. Combining these methods we \rev{calculated median ages and $68\%$ and $95\%$ confidence intervals} for the full sample. We are continuing to collect APO/ARCES spectra of the full sample over at least 3 epochs per star to identify stellar binaries (which could be the source of the astrometric acceleration).

We find approximately twice as many young stars from gyrochronology than expected from a uniform star formation rate. We identify two possible explanations for this result. First, there could be unidentified short-period binaries in the sample: either short period binaries that spin each other up, or visual binaries that overlap in \textit{TESS} images. Second, if binaries tend to dissolve over time, then young stars may be more likely to be part of binaries (and possess astrometric accelerations) than older stars. We expect our multi-epoch radial velocity monitoring to identify additional binaries in the sample, and we will report on the results of this monitoring in a future paper.

Looking ahead, \textit{Gaia} DR4 is expected to be released in mid-2026. Unlike DR3, which only contains a single position and velocity measurement, DR4 will contain individual scans over several years. These individual scans will better differentiate between short-period stellar binaries and long-period substellar companions. Future \textit{TESS} data releases will provide light curves for more accelerating stars, allowing for additional gyrochronological age measurements. Additional light curves for accelerating stars that already have light curves will allow for monitoring of spot evolution over the course of a stellar cycle.

The youngest of the accelerating stars in our sample are high-priority targets for direct imaging since these substellar companions will be bright enough to be imaged. In addition to continued radial velocity monitoring, we have begun a direct imaging campaign targeting the youngest stars. Results from both studies will be presented in future work. Substellar companions found in this sample will become key benchmarks for substellar evolution models.
